% Solutions to Chapter 10's exercises, from lectures/L10/appendix/e_solutions.md.
%
% Exercises 9 and 13 are programs: the programs themselves (9 b, 13 b) are left out, and what is
% worked on paper is kept: the flow chart in words, the table of test results and the
% explanations; the note on why 13 b) does the low nibble first follows 13 a). Exercise 3 c) keeps
% its four lines of assembly, since they are the answer asked for.
%
% One correction, made in the lecture too: 8 b) said that PB6 and PB7 are "kristallens stift", but
% the Uno clocks the ATmega328P with a ceramic resonator (section 11.1.7); it now says, as section
% 11.1.5 does, that the two pins are used by the clock.

\answerchapter{c10:ch}

Varje antal cykler här är uppmätt i simulatorn, inte bara uträknat.

\answer{c10:ex:1}{Tecken och koder}
\pt{a}\code{'5'} = \code{0x35}, \code{'9'} = \code{0x39}, \code{'E'} = \code{0x45}, \code{'F'} =
\code{0x46}.

\pt{b}\code{0x38 - 0x30 = 0x08}, alltså 8. Det är skälet till att \code{subi r16, 0x30} gör om ett
siffertecken till sitt värde.

\pt{c}Siffrorna \code{'0'}--\code{'9'} ligger direkt efter varandra från \code{0x30}, men mellan
\code{'9'} (\code{0x39}) och \code{'A'} (\code{0x41}) ligger sju andra tecken, och de måste hoppas
över.

\pt{d}För läsbarheten. \code{'E'} säger att konstanten är tänkt som ett tecken; \code{0x45} säger
bara att den är ett tal. Maskinen ser ingen skillnad, men den som läser programmet gör det.

\answer{c10:ex:2}{En byte blir text}
\pt{a}Den höga nibblen av \code{0x9E}:

\begin{booktable}{@{}llL@{}}
  \thead{Instruktion} & \thead{\code{r24}} & \thead{Kommentar} \\
  \midrule
  \code{mov r24, r16} & \code{0x9E} & \\
  \code{swap r24} & \code{0xE9} & halvorna byter plats \\
  \code{andi r24, 0x0F} & \code{0x09} & den höga nibblen \\
  \code{cpi r24, 10} & \code{0x09} & 9 < 10, så \code{brlo} hoppar över \code{subi r24, -7} \\
  \code{subi r24, -0x30} & \code{0x39} & \\
\end{booktable}

\noindent\code{r20} = \code{0x39} = \code{'9'}.

\pt{b}Den låga nibblen: \code{mov} ger \code{0x9E}, \code{andi} ger \code{0x0E}. 14 är inte mindre
än 10, så \code{brlo} hoppar inte, och \code{subi r24, -7} ger \code{0x15}. Sista \code{subi} ger
\code{0x45}. \code{r21} = \code{0x45} = \code{'E'}.

\pt{c}\code{0x39} på adress \code{0x0100} och \code{0x45} på \code{0x0101}: texten \code{9E}.

\pt{d}För \code{0x5A}: den höga nibblen blir \code{0x05}, som är mindre än 10, och blir
\code{0x35} = \code{'5'}. Den låga blir \code{0x0A}, som inte är mindre än 10:
\code{0x0A + 7 = 0x11}, och \code{0x11 + 0x30 = 0x41} = \code{'A'}. Texten blir \code{5A}.

\answer{c10:ex:3}{Från tecken till värde}
\pt{a}\code{0x42 - 0x30 = 0x12} = 18. 18 är inte mindre än 10, så koden subtraherar 7 till:
\code{r22} = 11 = \code{0x0B}. Rätt: B är 11.

\pt{b}\code{0x34 - 0x30 = 4}. 4 är mindre än 10, och koden är klar: \code{r22} = 4.

\pt{c}\code{0x62 - 0x30 = 0x32} = 50, och \code{50 - 7 = 43} = \code{0x2B}. Det är inget värde
0--15 alls; koden förutsätter stora bokstäver. En liten bokstav har en kod som är \code{0x20} högre
än motsvarande stora: \code{'a'} = \code{0x61} och \code{'A'} = \code{0x41}. Gör om små bokstäver
till stora först:

\needspace{6\baselineskip}
\begin{asmcode}
    cpi r22, 'a'                    ; A lowercase letter?
    brlo not_lower
    subi r22, 0x20                  ; 'a'-'f' -> 'A'-'F'.
not_lower:
\end{asmcode}

\noindent Det frestande alternativet, att nollställa bit~5 med \code{andi r22, 0xDF}, fungerar för
bokstäverna men förstör siffrorna: \code{'0'} = \code{0x30} har också bit~5 satt, och skulle bli
\code{0x10}.

\answer{c10:ex:4}{Registerpar}
\pt{a}\leavevmode

\begin{narrowtable}{@{}rccc@{}}
  \thead{Tal} & \thead{Hexadecimalt} & \thead{\code{high()}} & \thead{\code{low()}} \\
  \midrule
  2500 & \code{0x09C4} & \code{0x09} & \code{0xC4} \\
  256 & \code{0x0100} & \code{0x01} & \code{0x00} \\
  255 & \code{0x00FF} & \code{0x00} & \code{0xFF} \\
  40\,000 & \code{0x9C40} & \code{0x9C} & \code{0x40} \\
\end{narrowtable}

\noindent Uträkningen för 2500: \code{2500 / 256 = 9} med resten \code{196}, och
\code{196 = 0xC4}. För 40\,000: \mbox{\code{40 000 / 256 = 156}} med resten \code{64}, och
\code{156 = 0x9C}, \code{64 = 0x40}.

Lägg märke till 255 och 256: talen skiljer sig med ett, men bytena är helt olika. 255 är den
största låga byten, och 256 är den första ettan i den höga.

\pt{b}\code{0x12 · 256 + 0x34 = 18 · 256 + 52 = 4608 + 52 = 4660}.

\pt{c}\code{0xFFFF} = 65\,535 = \code{2^16 - 1}. Sexton bitar ger \code{2^16} = 65\,536 olika
värden, och det första är 0.

\answer{c10:ex:5}{Addition och subtraktion med 16 bitar}
\needspace{12\baselineskip}
\pt{a}\leavevmode

\begin{textdrawing}
                 r25    r24
                0x12   0xF0      4848
             +  0x0A   0x20      2592
             ---------------
  add r24:             0x110  -> r24 = 0x10, C = 1
  adc r25:     0x12 + 0x0A + 1 = 0x1D
             ---------------
                0x1D   0x10      7440
\end{textdrawing}

\noindent Kontroll decimalt: \code{0x12F0 = 18 · 256 + 240 = 4848},
\code{0x0A20 = 10 · 256 + 32 = 2592}, och \code{4848 + 2592 = 7440}. Resultatet
\code{0x1D10 = 29 · 256 + 16 = 7440}. Det stämmer.

\pt{b}\code{add} tar inte med minnessiffran: \code{0x12 + 0x0A = 0x1C}, och svaret blir
\code{0x1C10} = 7184. Det är 256 för lite, precis värdet av den minnessiffra som tappades.

\pt{c}\code{sub r20, r22}: \code{0x00 - 0x01} går inte utan att låna, så \code{r20} = \code{0xFF}
och \code{C} = 1. \code{sbc r21, r23}: \code{0x01 - 0x00 - 1 = 0x00}, så \code{r21} =
\code{0x00}. Resultatet är \code{0x00FF} = 255, som det ska.

\pt{d}\code{sbiw} fungerar bara på paren \code{r25:r24}, \code{r27:r26}, \code{r29:r28} och
\code{r31:r30} (\secref{c10:a8}), och assemblern godtar inte \code{r20}. Med talet i
\code{r25:r24} hade \code{sbiw r24, 1} fungerat.

\answer{c10:ex:6}{En fördröjning på 10 ms}
\pt{a}\mbox{\code{0,010 s · 16 000 000 cykler/s = 160 000 cykler}}.

\pt{b}De två \code{ldi} tar 2 cykler och loopen \code{4N - 1}, totalt \code{4N + 1}.
\mbox{\code{4N + 1 = 160 000}} ger \mbox{\code{N = 39 999,75}}, som inte är ett heltal.
\mbox{\code{N = 40 000}} ger \mbox{\code{160 001}} cykler, en cykel för mycket, och
\mbox{\code{N = 39 999}} ger \mbox{\code{159 997}}, tre för lite. \textbf{\mbox{\code{N = 40 000}},
160\,001 cykler}, vilket också är vad simulatorn mäter.

\pt{c}\code{r24} = 10, och anropet tar \mbox{\code{16 000 · 10 + 18 = 160 018}} cykler
(\secref{c10:c4}).

\pt{d}Loopen är närmast: en cykel, 62,5~ns, fel, mot 18 cykler, drygt 1~µs, för subrutinen. Men
båda felen är långt under vad något i ett styrsystem märker. Välj ändå subrutinen: den skrivs en
gång, tar tiden som parameter, och lämnar alla register som den fann dem. Loopen måste skrivas av
på varje ställe den behövs, och förstör \code{r25:r24} där.

\answer{c10:ex:7}{Knappen som läses baklänges}
\pt{a}När knappen är nedtryckt är stiftet direkt förbundet med jord, 0~V. Pull-up-motståndet på
20--50~kΩ ligger mellan 5~V och stiftet, och hela spänningen hamnar över motståndet, med en mycket
liten ström. Stiftet ligger på 0~V och läses som 0.

\pt{b}Utan pull-up är stiftet \textbf{flytande} när knappen är släppt: ingenting bestämmer dess
nivå. Programmet läser 0 eller 1 beroende på brus, luftfuktighet och var händerna är. Felet är
svårt att hitta eftersom programmet ofta \emph{verkar} fungera en stund, och slutar fungera när man
tar i kortet eller flyttar det.

\pt{c}En nedtryckt knapp läses som 0, men en lysdiod tänds av en etta. \code{com} vänder på varje
bit. Utan \code{com} hade alla lysdioder lyst när ingen knapp var nedtryckt, och en lysdiod hade
slocknat när dess knapp trycktes.

\answer{c10:ex:8}{Följ programmet}
\pt{a}Bit~2 är 0, så knappen på PD2 är nedtryckt. \code{com} ger \code{0b00000100}, och lysdioden
på PB2 lyser: \code{PORTB} = \code{0b00000100}.

\pt{b}Med \code{andi r16, 0xF0} behålls bara bit 4--7. För \code{0b11111011} blir det
\code{0b00000100 AND 0xF0} = \code{0b00000000}: ingenting lyser, trots att en knapp är nedtryckt,
eftersom den ligger i den bortmaskade halvan.

För \code{0b01101111}: \code{com} ger \code{0b10010000}, och masken ändrar ingenting.
\code{PORTB} = \code{0b10010000}, så PB7 och PB4 lyser. PB7 syns i simulatorn men inte på
Arduino-kortet, där PB6 och PB7 används av klockan.

\pt{c}\code{in} 1, \code{com} 1, \code{out} 1 och \code{rjmp} 2: 5 cykler per varv.
\mbox{\code{16 000 000 / 5 = 3 200 000}} läsningar per sekund.

\answer{c10:ex:9}{Start och stopp med självhållning}
\pt{a}Flödesplanen, i ord:
\begin{enumerate}
  \item Start. Gör PB0 till utgång, och slå på pull-uparna på PD2 och PD3.
  \item Är stoppknappen nedtryckt? \textbf{Ja:} släck PB0 och gå till 2. \textbf{Nej:} fortsätt.
  \item Är startknappen nedtryckt? \textbf{Ja:} tänd PB0. \textbf{Nej:} gör ingenting.
  \item Gå till 2.
\end{enumerate}
Frågan om stopp ställs \textbf{först}, så att stopp vinner när båda är nedtryckta.

\pt{c}Uppmätt i simulatorn:

\begin{narrowtable}{@{}lc@{}}
  \thead{Händelse} & \thead{\code{PORTB} bit~0} \\
  \midrule
  Ingen knapp & 0 \\
  Start trycks och släpps & 1, och förblir 1 \\
  Stopp trycks & 0 \\
  Båda nedtryckta & 0 \\
\end{narrowtable}

\pt{d}Minnet är utgångsregistret \code{PORTB} självt. När ingen knapp är nedtryckt gör programmet
ingenting med PB0, och biten behåller det värde den senast fick. Det motsvarar reläets
självhållningskontakt: den håller kvar tillståndet tills något aktivt ändrar det. Det är samma idé
som vippan i kapitel~\ref{c6:ch}, ett minne som ändras bara när det tillåts.

\answer{c10:ex:10}{\code{rcall}, \code{ret} och stacken}
\pt{a}Först läggs returadressen, adressen till instruktionen efter \code{rcall}, på stacken, och
\code{SP} flyttas två steg nedåt. Sedan hoppar programmet till subrutinen.

\pt{b}Returadressen måste finnas kvar tills subrutinen är klar, hur många register subrutinen än
använder och hur många anrop den själv gör. Ett fast register hade skrivits över av nästa anrop i
en subrutin som anropar en annan. Stacken ger varje pågående anrop en egen plats, och den växer och
krymper med anropen.

\pt{c}\code{0x08FD} i subrutinen, eftersom två byte har lagts dit. \code{0x08FF} igen efter
\code{ret}.

\pt{d}Allt som ligger under \code{SP} räknas som ledigt. Ingen läser det, och nästa \code{rcall}
eller \code{push} skriver över det.

\answer{c10:ex:11}{Kontroll: förutsäg stacken}
\pt{a}\code{rcall} ligger på \code{0x0007}, så returadressen är \code{0x0008}. \code{rcall} skriver
den låga byten först, på \code{0x08FF}, och den höga på \code{0x08FE}. Därefter lägger
\code{push r16} värdet \code{0xA5} på \code{0x08FD}, och \code{push r17} värdet \code{0x3C} på
\code{0x08FC}.

\begin{narrowtable}{@{}ccl@{}}
  \thead{Adress} & \thead{Byte} & \thead{Vad det är} \\
  \midrule
  \code{0x08FF} & \code{0x08} & returadressens låga byte \\
  \code{0x08FE} & \code{0x00} & returadressens höga byte \\
  \code{0x08FD} & \code{0xA5} & \code{r16} \\
  \code{0x08FC} & \code{0x3C} & \code{r17} \\
\end{narrowtable}

\noindent Fyra byte är upptagna, och \code{SP} = \code{0x08FF - 4} = \textbf{\code{0x08FB}}.

\pt{b}Simulatorn visar \code{SP} = \code{0x08FB}, och i fönstret Memory står, från adress
\code{0x08FC} och uppåt, bytena \code{3C A5 00 08}: precis tabellen ovan, läst nedifrån och upp.

\pt{c}De två vanligaste avvikelserna:
\begin{itemize}
  \item \textbf{Returadressen i fel ordning}, \code{0x00} på \code{0x08FF} och \code{0x08} på
    \code{0x08FE}. Den låga byten läggs dit först och hamnar därför på den högre adressen; stacken
    växer nedåt.
  \item \textbf{\code{SP} en byte fel}, \code{0x08FC}. \code{SP} pekar på den första
    \textbf{lediga} byten, inte på den senast använda. Den senast använda är \code{0x08FC}, och den
    första lediga är en under.
\end{itemize}

\pt{d}Efter \code{pop r17}, \code{pop r16} och \code{ret} är \code{SP} tillbaka på \code{0x08FF}.
Inget av de fyra bytena har ändrats: \code{pop} och \code{ret} läser och flyttar \code{SP}, men
raderar ingenting. \code{r16} och \code{r17} har sina gamla värden, \code{0xA5} och \code{0x3C}.

\answer{c10:ex:12}{Hitta felet}
\pt{a}Registren hämtas i samma ordning som de lades dit, i stället för i omvänd ordning.
\code{pop r16} får det som lades dit sist, \code{r17}:s värde, och tvärtom. Subrutinen återvänder
rätt, men programmet som anropar den finner \code{r16} och \code{r17} med bytta värden. Felet märks
först långt senare, som ett konstigt resultat, och ser inte alls ut att ha med subrutinen att göra.

\pt{b}En byte, \code{r16}:s sparade värde, ligger kvar ovanpå returadressen. \code{ret} läser den
som en del av adressen och hoppar någonstans där ingen har bestämt att programmet ska vara.
Därifrån kör processorn vad som helst som råkar ligga i programminnet, ofta tomt flashminne, tills
den kommer runt till adress 0 och börjar om. Det märks direkt, men felet syns inte på raden där det
sitter.

\pt{c}Läs \code{SP} i Processor Status när programmet står på \code{rcall}, och igen när det står
på instruktionen efter. Är värdena olika, lämnar subrutinen stacken i ett annat skick än den fann
den.

\answer{c10:ex:13}{Hela byten som text}
\pt{a}\code{r24} och \code{r25} får svaren och behöver inte sparas: det står i kommentaren att de
är subrutinens resultat. Subrutinen behöver också ett register för att spara byten under tiden,
här \code{r16}, och det måste sparas med \code{push} och \code{pop}.

Den låga nibblen görs först, eftersom \code{nibble_to_ascii} lämnar sitt svar i \code{r24}, och
det är där den höga nibblens tecken ska hamna till sist.

\pt{c}I simulatorn står \code{42 37} på adress \code{0x0100}, och fönstret visar \code{B7}. Stacken
är som djupast inne i \code{nibble_to_ascii}: två byte returadress till huvudprogrammet, en byte
för \code{r16} och två byte returadress till \code{byte_to_ascii}, fem byte. \code{SP} är då
\code{0x08FA}.
